%%%%%%%%%%%%%%%%%%%%%%%%%%%%%%%%%%%%%%%%%%%%%%%%%%%%%%%%%%%%%%%%%%%%%%%%%%%%%%%
% intro.tex: Introduction to the thesis
%%%%%%%%%%%%%%%%%%%%%%%%%%%%%%%%%%%%%%%%%%%%%%%%%%%%%%%%%%%%%%%%%%%%%%%%%%%%%%%%
\chapter{Introduction}
\label{chap:intro}
%%%%%%%%%%%%%%%%%%%%%%%%%%%%%%%%%%%%%%%%%%%%%%%%%%%%%%%%%%%%%%%%%%%%%%%%%%%%%%%%

The standard model (SM) of particle physics describes elementary particles and the forces by which they interact \cite{glashow_1961,weinberg,Salam:1968rm}. The SM is a phenomenological model built from over 100 years of experimental observation. It unifies the electromagnetic and weak forces and, along with the strong force, provides a coherent framework for the fundamental forces. Many successful predictions have been made using the SM, especially with respect to properties of particles which have been verified with unprecedented precision. One of these successful predictions is the Higgs mechanism and, by extension, the Higgs boson, which breaks the symmetry of electroweak interaction and explains how many particles obtain their mass~\cite{englert,higgs,higgs_1964,guralnic,higgs_1966,kibble_1967}. The Higgs boson was observed experimentally by the Compact Muon Solenoid (CMS) and A Toroidal LHC Apparatus (ATLAS) collaborations in 2012 using data collected at the European Center for Nuclear Research (CERN) Large Hadron Collider (LHC)~\cite{atlas:2012_higgs,cms:2012_higgs}. With the discovery of the Higgs boson, the particle content of the modern SM was completed.

While the SM does well to predict phenomena within its scope, there are many it does not describe at all. For example: the SM does not provide a description of the gravitational force, nor does it predict the existence of dark matter\footnote{A form of weakly interacting matter hypothesized as an explanation for some important phenomena in astrophysics} or dark energy\footnote{The hypothesized energy source responsible for the expansion of space-time.}~\cite{dark_matter_2020}. And even within its scope, the SM has incomplete descriptions of particles such as neutrinos. These particles are predicted to be massless, but experiments have confirmed that they are massive particles~\cite{Mertens_2016}. The SM also fails to address deeper problems in particle physics like the observed imbalance between matter and anti-matter in the universe~\cite{sakharov_1991} and the fine-tuning of the Higgs field vacuum expectation value~\cite{weinberg_1976,susskind_1979}. Finally, the SM is not an elegant or concise model; it has many free parameters that must be provided by experiment~\cite{review_of_particle_2020}. And yet, it is exactly this abundance of free parameters that this thesis targets. In fact, this shortcoming will be leveraged against itself. Making precision measurements of the SM's free parameters is a way to not only test its self-consistency but also to shed light on areas of the SM where increased precision magnifies differences between prediction and measurement. Any differences in prediction and measurement open the door to physics beyond the SM (BSM physics) and offer potential solutions to many of the problems listed above. 

Precision measurements of Higgs boson properties provide deeper insight into the theory of electroweak symmetry breaking\footnote{This is the mathematical framework used to describe the origins of the Higgs boson, as well as its interactions with the other gauge bosons.}, the nature of fermionic Yukawa interactions, and the shape of the Higgs field potential~\cite{ew_phase_transition_1996,ew_phase_transition_1999}. 
 
This thesis focuses on a precision measurement of the Higgs boson mass in the diphoton decay channel (H$\rightarrow\gamma\gamma$) using proton-proton collision data collected by the CMS experiment during Run 2 of the LHC. This thesis begins in Chapter~\ref{chap:theory} with a summary of the theoretical foundations of the SM that describe the Higgs boson and its relation to other particles, as well as summarizing the status of precision measurements of Higgs boson properties. Then, in Chapter~\ref{chap:cms}, the experimental apparatus, the CMS detector, is presented in detail and Chapter~\ref{chap:reco} describes how data collected by the CMS detector is processed into physics objects like electrons and photons to be used in analyses. Chapter~\ref{chap:scales} presents an improved method for correcting the energies of electrons and photons that allows analyses to make more precise measurements of particle properties. After this final set of corrections, Chapter~\ref{chap:categorisation} describes how H$\rightarrow\gamma\gamma$ events are grouped to improve the sensitivity of the Higgs boson mass measurement. Before discussion of the measurement itself, the statistical framework used to interpret results is laid out in Chapter~\ref{chap:stats}. Finally, the results of the Higgs boson mass measurement are presented in Chapter~\ref{chap:measurement} and a conclusion is given in Chapter~\ref{chap:summary}. 

%%%%%%%%%%%%%%%%%%%%%%%%%%%%%%%%%%%%%%%%%%%%%%%%%%%%%%%%%%%%%%%%%%%%%%%%%%%%%%%%